\subsection{GROUPS OPERATING CONTINUOUSLY ON A LOCALLY COMPACT SPACE}

PROPOSITION 9. Let G be a topological group operating continuously on a locally compact space X. Then if X/G is Hausdorff it is locally compact.

Since the equivalence relation on X defined by G is open (§ 2, no. 4, Lemma 2), the proposition results from Chapter I, § 10, no. 4, Proposition 10.

PROPOSITION 10. Let G be a topological group operating continuously on a locally compact space X, and suppose that X/G is Hausdorff. Let \(\varphi\) be the canonical mapping of X onto X/G. Then if K' is any compact subset of X/G, there is a compact subset K of X such that \(\varphi(K) = K'\) .

Since the equivalence relation defined by G is open (§ 2, no. 4, Lemma 2), the proposition is a particular case of Proposition 10 of Chapter I, § 10, no. 4.

PROPOSITION 11. Let G be a Hausdorff topological group operating properly on a non-empty space X. If X is compact (resp. locally compact) then so are G and X/G.

By hypothesis, the mapping \(\theta : (s, x) \to (x, s.x)\) of \(G \times X\) into \(X \times X\) is proper; if \(X \times X\) is compact (resp. locally compact) then the Corollary to Proposition 9 of Chapter I, § 10, no. 4 shows that \(G \times X\) is also compact (resp. locally compact), and therefore so is \(G\) since \(X \neq \emptyset\) . Since \(X/G\) is Hausdorff (no. 2, Proposition 3), the compactness (resp. local compactness) of \(X\) implies the compactness (resp. local compactness) of \(X/G\) {[}Chapter I, § 10, no. 4, Proposition 8 (resp. Proposition 9){]} (see § 2, Exercise 29).

We shall now give criteria which allow us to assert that a Hausdorff topological group G operates properly on a locally compact space X. For each pair of subsets K, L of X we denote by P(K, L) the set of all \(s \in G\) such that \(s \cdot K \cap L \neq \emptyset\) .

THEOREM 1. Let G be a Hausdorff topological group which operates continuously on a topological space X. Let K be a compact subset of X and let L be a closed subset of X. Then:

\begin{enumerate}
\def\labelenumi{\alph{enumi})}
\item
  The set \(\mathbf{P}(\mathbf{K},\mathbf{L})\) is closed in \(\mathbf{G}\) .
\item
  If G operates properly on X and if L is compact, then P(K, L) is compact.
\item
  Conversely, if X is locally compact and if, for each pair K, L of compact subsets of X, P(K, L) is relatively compact in G {[}and therefore compact by a){]}; then G operates properly on X (and if X is not empty, G is locally compact by Proposition 11).
\end{enumerate}

The mapping \((s, x) \to s.x\) of \(G \times K\) into X is continuous, and the inverse image \(L'\) of L under this mapping is therefore closed. Since K is compact, the projection \(pr_{1}: G \times K \to G\) is proper (Chapter I, § 10, no. 2, Corollary 5 of Theorem 1) and the image of \(L'\) under \(pr_{1}\) is therefore closed. This image is \(P(K, L)\) , hence a) is proved.

To prove b): X is Hausdorff (no. 2, Proposition 3). By hypothesis the mapping \(\theta: (s, x) \to (x, s.x)\) of \(G \times X\) into \(X \times X\) is proper; since \(K \times L\) is compact, so is \(\theta^{-1}(K \times L)\) since \(X\) is Hausdorff (Chapter I, § 10, no. 2, Proposition 6). Hence the projection \(P(K, L)\) of \(\theta^{-1}(K \times L)\) into \(G\) is a compact set.

To prove c): since \(K \times L\) is closed in \(X \times X\) , it follows that \(\theta^{-1}(K \times L)\) is closed in \(P(K, L) \times K\) and is therefore compact under the hypotheses of c). Since every compact subset of \(X \times X\) is contained in a compact subset of the form \(K \times L\) , it follows that the inverse image of any compact subset of \(X \times X\) under \(\theta\) is compact, and since \(X \times X\) is locally compact, this shows that \(\theta\) is proper (Chapter I, § 10, no. 3, Proposition 7) {[}see Chapter IV, § 1, Exercise 4c){]}.

Remark. Clearly, we have \(\mathrm{P}(\mathrm{K},\mathrm{L})\subset\mathrm{P}(\mathrm{K}\cup\mathrm{L},\mathrm{K}\cup\mathrm{L})\) ; therefore for G to operate properly on a locally compact space X, it is sufficient that, for each compact subset K of X, the set \(\mathrm{P}(\mathrm{K},\mathrm{K})\) is relatively compact in G. In particular, a discrete group G operates properly on a locally compact space X if and only if, for each compact subset K of X, the set of \(s\in G\) such that \(s.K\cap K\neq\emptyset\) is finite.

Example. * Let X be a complex analytic manifold, analytically isomorphic to a bounded open subset of \(\mathbf{C}^n\) , and let G be the group of analytic automorphisms of X. The topology of compact convergence is compatible with the group structure of G, and it can be shown that G operates properly on X. In particular, every discrete subgroup of G operates properly on X.

Take for example X to be the upper half-plane \(\mathfrak{J}(z)>0\) , which is analytically isomorphic to an open disc in C. Then G is the group of all transformations \(z\to (az+b)/(cz+d)\) , where \(a,b,c,d\) are real and \(ad-bc\neq0\) . The subgroup H of G which consists of all such transformations for which \(a,b,c,d\) are integers and \(ad-bc=1\) is a discrete subgroup of G, called the modular group. By what has been said above, it operates properly on the upper half-plane \(\mathfrak{J}(z)>0\) . *

PROPOSITION 12. Let G be a Hausdorff topological group which operates continuously on a topological space X. Let K be a compact subset of X, and let \(\rho_{K}\) be the mapping \((s, x) \to s.x\) of \(G \times K\) into X. Then:

\begin{enumerate}
\def\labelenumi{\alph{enumi})}
\item
  If G operates properly on X, \(\rho_{\mathbf{K}}\) is proper.
\item
  If X is locally compact and \(\rho_{K}\) is proper for each compact subset K of X, G operates properly on X.
\end{enumerate}

The mapping \(\rho_{\mathbf{K}}\) factorizes as \(\mathbf{G} \times \mathbf{K} \xrightarrow{\theta_{\mathbf{K}}} \mathbf{K} \times \mathbf{X} \xrightarrow{\mathrm{pr}_2} \mathbf{X}\) , where \(\theta_{\mathbf{K}}\) is the restriction to \(\mathbf{G} \times \mathbf{K}\) of the mapping \(\theta: (s, x) \to (x, s.x)\) of \(\mathbf{G} \times \mathbf{X}\) into \(\mathbf{X} \times \mathbf{X}\) . Since \(\theta^{-1}(\mathbf{K} \times \mathbf{X}) = \mathbf{G} \times \mathbf{K}\) , \(\theta_{\mathbf{K}}\) is proper if \(\theta\) is proper (Chapter I, § 10, no. 1, Proposition 3). On the other hand, since \(\mathbf{K}\) is compact, the projection \(\mathrm{pr}_2: \mathbf{K} \times \mathbf{X} \to \mathbf{X}\) is proper (Chapter I, § 10, no. 2, Theorem 1, Corollary 5), hence \(\rho_{\mathbf{K}}\) is proper (ibid., no. 1, Proposition 5).

Suppose conversely that \(\rho_{\mathbf{K}}\) is proper for each compact \(\mathbf{K} \subset \mathbf{X}\) . If \(\mathbf{L}\) is a compact subset of \(\mathbf{X}\) , then \(\rho_{\mathbf{K}}^{-1}(\mathbf{L})\) is a compact subset of \(\mathbf{G} \times \mathbf{K}\) , whose projection into \(\mathbf{G}\) is \(\mathrm{P}(\mathbf{K}, \mathbf{L})\) ; therefore \(\mathrm{P}(\mathbf{K}, \mathbf{L})\) is compact. Hence if \(\mathbf{X}\) is locally compact it follows from Theorem 1 that \(\mathbf{G}\) operates properly on \(\mathbf{X}\) .

COROLLARY. Let G be a Hausdorff topological group operating properly on a topological space X. If K is any compact subset of X and if F is any closed subset of G, then F.K is closed in X.

This follows from Proposition 12 and from Chapter I, § 10, no. 1, Proposition 1.

